\section{Methods}\label{methods}

A world in SCORE is generated once, offline, in cloud batches, from the creator's picks for each place: references, a concept picture and a style. One world spans the scales its story needs (the game world of 2099 holds rooms, the ground of the Moon and the ground of Mars); scales far from those, such as galaxies or cells, make worlds of their own with their own solvers. The output is one canonical scene, an OpenUSD stage \citep{openusd} with glTF geometry \citep{gltf2}, in which every object carries its kind, place in metres, collision, the surface of each part, its sound, and tags for what a person can do with it, such as door, seat or airlock. Engines and simulators load it through thin adapters. The framework writes a generated base layer; the creator's own changes live in edit layers above it.

Between picks and scene, coding agents and open models work through a fixed sequence of stages: a dimensioned plan in metres, an inventory with a row for every object, one clean close-up picture per object, a shape for each object, its surfaces, its sounds, and assembly. Each stage's output is kept. The creator can therefore trace a weak result to the stage that caused it (a bad close-up, a bad 3D shape, a bad surface), rerun or hand-fix that stage alone, and regenerate everything downstream; the regeneration replaces the base layer and keeps every edit.

Every stage that costs money is preceded by deterministic checks on real geometry rather than an agent's judgement, because a fault costs nothing in the plan, cents at a picture, tens of cents and an hour at a 3D model, and an evening of the creator's time once a place is built. The plan is checked for leaks, blocked openings, pieces off their surface and doors that do not separate their two sides; the inventory against every element of the concept; each model for being closed, thick enough and straight; the surfaces against the place's palette and budget; and the assembled scene for anything floating or sinking and any visible mesh the route did not make (Appendix D). A failure sends the work back one stage.

The shape stage decides, per kind of object, between a code builder written by a coding agent and a generated model. Picture-to-3D models build chunky solids well and flat panels, thin beams and openings badly; code builders are exact and straight but show only the parts the agent wrote. The \textbf{parts check} settles it: a kind is built in code only when its build shows every part its close-up shows, and is otherwise generated from the close-up with Pixal3D \citep{li2026pixal3d} and made solid. Composites are split in the inventory, one model per real-world object, so a console is a desk carrying its monitors and keyboards rather than one fused mesh.

Shapes carry no surface of their own. Every part takes one surface from a shared library of ProcFunc functions \citep{raistrick2026procfunc} built on Infinigen's shaders \citep{raistrick2023infinigen}, coloured from the place's palette, with wear, dirt and seed as settings and one wear setting per room, applied from causes such as edges and foot traffic; parts of a generated model come from PartCrafter \citep{lin2025partcrafter}, and its picture only chooses which library surface each part takes. Sound follows the same pattern: surfaces carry their footsteps and impacts, rooms their echo and objects their own sound, generated by MOSS-SoundEffect \citep{moss2026} and chosen against its prompt by CLAP \citep{wu2022clap} (Appendix F). Both libraries grow with each world and are reused in the next.

The creator's review comes last. SCORE supplies the tools for it: pages with each stage's outputs side by side, before and after shots from the same cameras, in-engine walkthroughs, and every check's result; how the creator judges with them is the creator's own.

\begin{center}\fbox{\begin{minipage}{0.94\linewidth}\setlength{\parskip}{4pt}\small

\textbf{Gap: Figure 2, the framework.} The stages with their kept outputs, the scene's base and edit layers, and the engine adapters. To be drawn with the scene export, which is being built in the framework repository.

\end{minipage}}\end{center}

\begin{center}\fbox{\begin{minipage}{0.94\linewidth}\setlength{\parskip}{4pt}\small

\textbf{Gap: Figure 3, one place through the stages.} One place from concept to finished world, with each stage's output and the checks that fired on it. Waits on a place built and accepted on the final route.

\end{minipage}}\end{center}
